% =========================================================================
% Width-Bounded Heuristics for Symbolic Search
% Draft for ICAPS 2027. Compile with pdflatex; port to icaps27.sty later.
% Experiments section reports results from the SymK implementation.
% =========================================================================
\documentclass[letterpaper]{article}
\usepackage[margin=1in]{geometry}
\usepackage{amsmath,amssymb,amsthm}
\usepackage{natbib}
\usepackage{booktabs}
\usepackage{microtype}

% ------------------------------- macros ---------------------------------
\newcommand{\inlinecite}[1]{\citet{#1}}
\newcommand{\task}{\ensuremath{\Pi}}
\newcommand{\vars}{\ensuremath{\mathcal{V}}}
\newcommand{\ops}{\ensuremath{\mathcal{O}}}
\newcommand{\states}{\ensuremath{\mathcal{S}}}
\newcommand{\sinit}{\ensuremath{s_{\mathrm{I}}}}
\newcommand{\sgoal}{\ensuremath{s_\star}}
\newcommand{\pre}{\ensuremath{\mathit{pre}}}
\newcommand{\eff}{\ensuremath{\mathit{eff}}}
\newcommand{\cost}{\ensuremath{\mathit{cost}}}
\newcommand{\hstar}{\ensuremath{h^{*}}}
\newcommand{\hblind}{\ensuremath{h^{0}}}
\newcommand{\hpot}{\ensuremath{h^{\text{pot}}}}
\newcommand{\gstar}{\ensuremath{g^{*}}}
\newcommand{\copt}{\ensuremath{C^{*}}}
\newcommand{\bdd}[1]{\ensuremath{B_{#1}}}
\newcommand{\bddsize}[1]{\ensuremath{|B_{#1}|}}
\newcommand{\cofcount}{\ensuremath{Q}}
\newcommand{\width}{\ensuremath{\mathsf{w}}}
\newcommand{\values}{\ensuremath{V}}
\newcommand{\level}[2]{\ensuremath{H_{#2}}}
\newcommand{\effort}{\ensuremath{\mathit{effort}}}
\newcommand{\sub}[2]{\ensuremath{\mathrm{cof}_{#1}(#2)}}
\newcommand{\img}{\ensuremath{\mathrm{img}}}
\newcommand{\bddastar}{BDDA\ensuremath{^{*}}}
\newcommand{\astar}{A\ensuremath{^{*}}}
\newcommand{\ghseta}{GHSETA\ensuremath{^{*}}}
\newcommand{\symba}{SymBA\ensuremath{^{*}}}

\theoremstyle{plain}
\newtheorem{theorem}{Theorem}
\newtheorem{lemma}{Lemma}
\newtheorem{proposition}{Proposition}
\newtheorem{corollary}{Corollary}
\theoremstyle{definition}
\newtheorem{definition}{Definition}
\newtheorem{example}{Example}
\theoremstyle{remark}
\newtheorem{remark}{Remark}

\title{The Cost of Guidance in Symbolic Search}
\author{Anonymous for review}
\date{}

\newtheorem{fact}{Fact}
\begin{document}
\maketitle

\begin{abstract}
Heuristic symbolic search has been analyzed through a single lens: the
cumulative BDD size of the expanded state sets. Under this measure,
heuristics of bounded cofactor width are provably safe---partitioning
frontiers by their values costs at most a small polynomial factor---yet
in practice such heuristics still lose coverage against blind search.
We resolve this paradox by refining the cost model: expanding each of up
to $\values$ heuristic-value buckets per layer triggers its own image
operation, an operational cost that node counts cannot see and that no
width bound caps. The refined model makes a testable prediction---%
$f$-ordered expansion pays iff the number of heuristic values is
small---and suggests an algorithm: \emph{pruning-only} search keeps the
single-BDD layer structure of blind search, uses the heuristic solely
for provably sound discards, matches blind search's image-call count,
and improves the known effort bound by a factor of the width.
Experiments on the IPC benchmarks confirm both: merge-and-shrink
heuristics (few values) gain coverage through ordering while potential
heuristics (many values) lose it and are repaired by pruning; and a
four-design study of bidirectional search---culminating in seeding an
incumbent from a satisficing planner---shows that the binding constraint
at the state of the art is not fragmentation but the cost of computing
the heuristic at all, while the incumbent itself emerges as the cheapest
useful form of guidance, improving even blind search.
\end{abstract}

\section{Introduction}

Symbolic search with binary decision diagrams
\citep[BDDs;][]{bryant-ieeecomp1986,mcmillan-1993,edelkamp-helmert-ecp1999}
is among the strongest approaches to cost-optimal classical planning,
and its strongest configurations are blind
\citep{torralba-et-al-aij2017,speck-et-al-jair2025}. The classical
explanation is representational: partitioning a frontier BDD into
buckets by heuristic value, as \bddastar\ \citep{edelkamp-reffel-ki1998}
does, can inflate BDD sizes---exponentially in the worst case, even for
the perfect heuristic \citep{speck-et-al-icaps2020}. A recent
counterpoint shows the inflation is bounded by a polynomial in the
heuristic's \emph{cofactor width}, the maximal number of distinct
cofactors of the heuristic function per level of the search's variable
order, and that standard heuristic families have small width (see the
companion paper, cited here as \emph{Widths}\footnote{Reference
suppressed for double-blind review; the companion manuscript is included
as supplementary material. Its results used here are restated as
Facts~\ref{thm-bucket}--\ref{thm-effort} below.}). And yet, measured on
the IPC benchmarks, width-bounded heuristics with bounded fragmentation
\emph{still} lose coverage against blind forward search. Safety, it
turns out, is not the whole story.

This paper locates the missing cost. The effort measure used throughout
this literature \citep{speck-et-al-icaps2020} counts the BDD nodes of
the expanded sets; it does not count \emph{image operations}. Expanding
a layer bucket by bucket multiplies the number of image computations by
the number of nonempty heuristic values $\values$, and each image
operation carries overhead beyond the size of its argument. We make this
operational cost explicit in an image-aware effort measure, prove that
no width bound can cap the resulting factor for \bddastar, and derive
two consequences. First, a prediction: $f$-ordered expansion should pay
off precisely when $\values$ is small. Second, an algorithm:
\emph{pruning-only search} keeps every layer in a single BDD exactly
like blind search, uses an admissible heuristic only to discard states
that provably cannot improve on the incumbent solution (dead ends, and
the slice $g + h \geq c$ under upper bound $c$), performs one image
operation per layer, and---via an exact correspondence between its
layers and value slices of the blind layers---carries an effort bound a
factor $W$ tighter than that of \bddastar.

Our experiments in SymK \citep{speck-et-al-jair2025} confirm the
prediction and probe the limits of the algorithm. Linear merge-and-shrink
heuristics, with a median of $16$ values, convert ordering into a
coverage gain over blind forward search; potential heuristics
\citep{pommerening-et-al-aaai2015,fiser-et-al-aij2024}, with a median of
$53$, lose coverage through ordering and are restored to parity by
pruning. A four-design study of bidirectional search---eager, budgeted,
deferred, and finally seeded with an incumbent from a satisficing
planner---shows a consistent picture: pruning is safe and its search is
free, but the cost of \emph{computing} the heuristic is not reliably
recouped against blind bidirectional search, whereas the incumbent
itself improves even the blind configuration. The emerging design rules
are simple: order only with few-valued heuristics, prune with the rest,
acquire an incumbent as early as possible, and spend on heuristic
construction only what the proof phase can pay back.

\section{Background}

We consider cost-optimal planning for SAS$^{+}$ tasks
\citep{backstrom-nebel-compint1995} with positive integer operator
costs, binary-encoded into $n$ Boolean variables under a fixed variable
order $\pi$ as usual in symbolic search
\citep{torralba-et-al-aij2017,kissmann-hoffmann-jair2014}; $|B|$ denotes
the inner-node count of a BDD $B$. Symbolic uniform-cost (blind) search
expands, per $g$-value, the layer $L_g$ of states first reached at cost
$g$, computing one image per layer; \bddastar\
\citep{edelkamp-reffel-ki1998} splits each layer into buckets
$L_g \cap h^{-1}(v)$ and expands them in order of increasing $f = g+v$,
ties broken by smaller $g$. The \emph{effort} of a run is the cumulative
size of the BDDs expanded with $g + v \leq \copt$ and $g < \copt$
\citep{speck-et-al-icaps2020}, where \copt\ is the optimal cost.

The \emph{cofactor width} $\width_\pi(h)$ of a heuristic $h$ is the
maximal number of distinct cofactors of $h$ (as a function over the
Boolean encoding) after any prefix of $\pi$---equivalently, the maximal
number of nodes per level of its quasi-reduced algebraic decision
diagram \citep{wegener-2000}. We write $\values \leq \width_\pi(h)$ for
the number of finite heuristic values. Two results of the companion
paper are used throughout; we restate them without proof.

\begin{fact}[bucket bound, \emph{Widths}]
\label{thm-bucket}
For every state set $S$, every $U \subseteq \img(h)$, and
$W = \width_\pi(h)$:
$|B_{S \cap h^{-1}(U)}| \leq W n \, (|B_S| + 2)$.
\end{fact}

\begin{fact}[effort bound, \emph{Widths}]
\label{thm-effort}
For consistent $h$ with $\width_\pi(h) \leq W$, the effort of \bddastar\
is at most $W^2 n \cdot (\mathit{effort}(\hblind) + 2\copt)$, and each
expanded bucket with key $(g,v)$ is exactly $L_g \cap h^{-1}(v)$
(bucket identity, \emph{Widths}).
\end{fact}

We also reuse the companion paper's width-bounded families: integer
potential heuristics with magnitude cap $M$
\citep{fiser-et-al-aij2024}, prefix pattern databases, and linear
merge-and-shrink abstractions \citep{helmert-et-al-jacm2014} with shrink
limit $N$, whose level sets $H_v = h^{-1}(v)$ we compute as BDDs in the
search order. For the moved definitions referenced below we write
$\effort(\task, \hblind)$ for blind forward effort and
$N_\ell$ for the number of nonempty layers.

\section{A Refined Cost Model: Image Calls}
\label{sec-cost}

The effort measure inherits from \citet{speck-et-al-icaps2020} a
cost model that counts the BDD nodes of the expanded sets. This captures
the representational cost of partitioning, but not its \emph{operational}
cost: every expanded bucket triggers its own image computation, and an
image operation carries overhead beyond the size of its argument---the
traversal of the transition relations and the associated cache
effects. We make this explicit.

\begin{definition}[image-aware effort]
\label{def-image-effort}
For $\alpha \geq 0$, the \emph{$\alpha$-effort} of a symbolic search run
is $\effort_\alpha = \sum_{B} (\alpha + |B|)$, where the sum ranges over
the image operations performed before termination and $B$ is the
argument of each. The parameter $\alpha$ abstracts the per-call
overhead; for $\alpha = 0$ and the expansions counted by
the standard effort measure, the two coincide.
\end{definition}

\begin{proposition}[image-call counts]
\label{prop-calls}
On a task with $N_\ell$ nonempty layers, blind forward search performs
$N_\ell$ image operations. \bddastar\ with heuristic $h$ performs
$\sum_{g} \values_g$ image operations, where $\values_g \leq \values$ is
the number of nonempty buckets of layer $g$; the total can thus exceed
the blind count by a factor of up to $\values$, the number of finite
heuristic values, regardless of the width and of the total BDD size.
\end{proposition}

\begin{proof}
Immediate from the algorithms: blind search computes one image per
expanded layer, and \bddastar\ expands every nonempty $(g, v)$ bucket
separately (the bucket identity, Fact~\ref{thm-effort}), computing one image each; layer
$g$ has $\values_g \leq |\img(h)| = \values$ nonempty buckets.
\end{proof}

Fact~\ref{thm-effort} bounds the node term of $\effort_\alpha$ for
every $\alpha$, but no width bound can cap the $\alpha$-term of \bddastar:
the factor $\values_g$ is inherent to expanding buckets separately. This
gap between the two cost models is not hypothetical---in our experiments
(Section~\ref{sec-experiments}), the tasks where \bddastar\ with
potentials loses coverage against blind search are exactly the tasks with
many distinct heuristic values (median $65$), while the node-based
fragmentation stays small everywhere. The remedy is an algorithm whose
image-call count matches blind search by construction.

\section{Pruning-Only Search}
\label{sec-prune}

The preceding analysis motivates a
variant that keeps every layer in a \emph{single} BDD, exactly like blind
forward search, and uses the heuristic only to discard states that
provably cannot contribute a better plan. \emph{Pruning-only search} runs
symbolic uniform-cost search and maintains the incumbent bound $c$, the
cost of the best solution found so far ($c = \infty$ initially; symbolic
forward search discovers solutions through frontier--goal cuts before
optimality is proved). From each generated layer it discards the states
with $h(s) = \infty$ and the slice $\{\, s \mid g + h(s) \geq c \,\}$,
re-applying the current bound when the layer is expanded. Admissibility
of $h$ makes both discards sound.

\begin{lemma}[slice identity]
\label{lem-slice}
Let $h$ be consistent, and let $c_g$ denote the incumbent bound when
pruning-only search expands the layer with $g$-value $g$. Since layers
are expanded in order of increasing $g$ and the incumbent bound never
increases, the expanded layer is exactly
$L_g \cap h^{-1}(\{\, v \mid v < c_g - g \,\})$.
\end{lemma}

\begin{proof}
The expanded set is contained in the slice because the bound $c_g$ is
applied at expansion time. For the converse, let
$s \in L_g$ with $g + h(s) < c_g$ and let
$p_0 = \sinit, \dots, p_k = s$ be a cheapest path, so
$\gstar(p_j) + h(p_j) \leq \gstar(s) + h(s) < c_g$ for every $j$ by
consistency along the path. Layer $\gstar(p_j)$ was expanded no later
than layer $g$, so its bound satisfied $c_{\gstar(p_j)} \geq c_g$, and
$p_j$ survived both its generation-time and expansion-time pruning; by
induction along the path, $s$ is generated at cost $g$ and survives, so
it lies in the expanded layer.
\end{proof}

Each expanded layer is thus a single value slice of the blind layer, so
Fact~\ref{thm-bucket} with $U = \{\, v \mid v < c_g - g \,\}$ applies
directly, and only one such set is expanded per layer:

\begin{corollary}[pruning bound]
\label{cor-prune}
Let $h$ be admissible and consistent with $\width_\pi(h) \leq W$. The
cumulative size of the layers expanded by pruning-only search is at most
$W n \cdot (\effort(\task, \hblind) + 2\copt)$, one factor $W$ less than
the \bddastar\ bound of Fact~\ref{thm-effort}, with one image
operation per layer as in blind search. With $c = \infty$ throughout
(no solution found yet) and no dead ends, the search coincides with
blind forward search.
\end{corollary}

Pruning-only search trades the $f$-ordered expansion of \bddastar\ for
the layer structure of blind search: it reaches the first solution at
the same time as blind search, but then prunes the remaining
optimality-proof effort by the $g + h \geq c$ slices, and it discards
dead ends from the start. In the refined cost model its
$\alpha$-effort is at most
$W n \cdot (\effort(\task, \hblind) + 2\copt) + \alpha N_\ell$ for every
$\alpha$---the same image-call term as blind search, in contrast to
Proposition~\ref{prop-calls}. It thereby turns the width bound into a
pure safety guarantee: never more than a factor $Wn$ worse than blind
search, never an extra image operation, and any pruning is a strict
saving.

Unlike the $f$-ordered expansion of \bddastar, whose analysis is tied to
forward search, the pruning argument is direction-agnostic, so it
extends to the practically strongest configuration, bidirectional blind
search. Each direction uses its own admissible estimate: the forward
frontier discards $\{\, s \mid g_{\mathrm{fw}} + h(s) \geq c \,\}$ with
$h$ estimating goal distance, the backward frontier
$\{\, s \mid g_{\mathrm{bw}} + \bar{h}(s) \geq c \,\}$ with $\bar{h}$
estimating the distance \emph{from the initial state}. Fixing the
direction schedule and the sequence of incumbent bounds,
Lemma~\ref{lem-slice} applies within each direction, so the cumulative
expanded size is at most a factor
$\max(\width_\pi(h), \width_\pi(\bar h)) \cdot n$ times that of blind
bidirectional search under the same schedule, with no additional image
operations. We state this as an observation rather than a theorem:
pruning can shift \emph{when} solutions are discovered, which feeds back
into the bound sequence and, in adaptive implementations, into the
schedule itself; a self-contained formal treatment must account for this
feedback and is left for future work. A single linear merge-and-shrink abstraction
conveniently supplies both estimates---the goal distances and the init
distances of the same shrunk transition system---and, because pruning
never partitions, this variant also requires no restriction to positive
operator costs for its soundness.


\section{Experiments}
\label{sec-experiments}

The setup follows the companion paper: all IPC 1998--2023 optimal-track
tasks without conditional effects and axioms; SymK on Intel Xeon Gold
6130 processors, 30 minutes and 8 GiB per run, Downward Lab
\citep{seipp-et-al-zenodo2017}. Because the heuristic searches assume
positive operator costs, forward comparisons use the 1377 positive-cost
tasks; the bidirectional study uses all 1697 tasks (pruning-only search
needs no cost restriction), and runs are always compared within the same
sweep. Coverage counts proved-optimal outcomes only.

\begin{table}[t]
\centering
\begin{tabular}{lrr}
\toprule
Configuration & \bddastar & pruning-only \\
\midrule
Blind forward (SymK) & \multicolumn{2}{c}{682} \\
Potentials ($M=8$) & 622 & 683 \\
Prefix PDB & 685 & 682 \\
Linear merge-and-shrink & 704 & 681 \\
\bottomrule
\end{tabular}
\caption{Coverage on the 1377 positive-cost tasks: $f$-ordered
(\bddastar) versus pruning-only use of each width-bounded family.}
\label{tab-coverage}
\end{table}

\paragraph{Ordering versus pruning.}
Table~\ref{tab-coverage} also separates the two ways of using a
heuristic. The pruning-only variants (Corollary~\ref{cor-prune}) behave
exactly as the theory promises: their coverage matches blind forward
search within one task for all three families ($683$, $682$, $681$
vs.\ $682$), their expansion-size ratio over blind search is $1.00$ in
the geometric mean, and where they deviate they save effort (dead-end
pruning on up to $20$ tasks per family). Runtime tells the same story
with one nuance the node-based measure misses: the geometric-mean
runtime ratio over blind search is $1.02$ for pruning with potentials
but $1.20$ and $1.44$ for prefix PDBs and merge-and-shrink, whose
heuristic \emph{construction} time is included---the search itself is
free, computing the heuristic is not. Pruning also gains little in
forward search, because the anytime upper bound only becomes available
late. Whether the $f$-ordered expansion of
\bddastar\ pays instead is decided by the image-call factor of
Proposition~\ref{prop-calls}: linear merge-and-shrink, whose level count
is small (median $16$ values), profits from ordering ($704$ vs.\ $681$
pruning-only, at a runtime ratio of $1.42$ that its extra coverage
justifies), while potentials, with a median of $53$ values and $65$
where losses occur, are hurt by it ($622$ vs.\ $683$; runtime ratio
$1.77$). The choice between guidance and safety is thus governed by the
number of heuristic values, exactly as the refined cost model predicts.


\paragraph{Bidirectional pruning.}
Since pruning needs no partitioning, it extends to blind bidirectional
search (Section~\ref{sec-prune}), whose coverage no forward
configuration approaches. Here the abstraction itself becomes the cost:
without a construction budget, building the merge-and-shrink level sets
consumed up to the entire time limit and cost more tasks than the
pruning recovered ($1032$ vs.\ $1051$ for blind bidirectional search on
the full suite of 1697 tasks, with construction taking over $90$ seconds
on ten percent of the tasks; a tenfold larger abstraction only widens
the gap, to $990$). With a $60$-second construction budget and
fallback to blind search---making the pruning safe \emph{including} its
setup cost---the budgeted variant recovers most of the difference
($1041$ vs.\ $1051$). Deferring construction further, to the moment the
first solution makes the bound slice applicable---so that the search is
identical to blind bidirectional search up to that point---narrows the
gap again ($1046$ vs.\ $1051$; one win, six losses), but does not close
it: on the losing tasks the optimality proof consumes nearly the whole
remaining budget, and the sixty seconds invested in the abstraction
prune too little to pay for themselves. Across all three designs the
picture is consistent, and we report it as an instructive negative
result: bidirectional blind search is robust not because heuristic
information fragments its BDDs---our bounds and measurements rule that
out---but because at this level of performance even provably safe
pruning must first pay for computing its heuristic, and on this suite
that constant cost is not reliably recouped. The potential is
nevertheless real: the virtual best of blind bidirectional search and
the forward merge-and-shrink configuration solves $1081$ tasks---the
forward search solves $30$ tasks the bidirectional one cannot---but the
unique solves are slow (median $947$ seconds), so neither time-slicing
nor simple pre-search selection harvests the complementarity; doing so
is an open problem our data makes concrete.


\paragraph{Bound seeding.}
Since the $g + h \geq c$ slice can only prune once an incumbent exists,
we finally supplied one externally: a $60$-second run of the satisficing
planner \textsc{LAMA} yields a plan of cost $c$, and the symbolic search
then runs with the exclusive bound $c$---it either finds a cheaper plan
or, by exhausting the bounded space, proves the incumbent optimal. All
configurations receive the same total budget, with seeding time charged
to the seeded ones, and coverage counts only proved-optimal outcomes.
The results reorder the picture in an unexpected way. Seeding helps
plain blind bidirectional search itself ($1054$ vs.\ $1050$): with a
bound it never needs to find or reconstruct a solution, only to raise
its lower bound, and this saving alone exceeds the seeding cost. Adding
merge-and-shrink pruning on top gains nothing ($1051$; it trails its own
seedless-pruning control)---the construction cost again consumes the
pruning gain. In the forward direction, by contrast, seeding completes
the pruning story positively: pruning-only potentials reach $688$
positive-cost tasks against $682$ for blind forward search and $683$
unseeded, the first forward-potential configuration to show a net gain.
The lesson is symmetric to the bidirectional one: an incumbent is the
cheapest form of guidance---it even helps blind search---and heuristic
pruning pays off on top of it exactly when the heuristic costs
essentially nothing to build.


\section{Related Work}

The effort measure we refine originates with
\citet{speck-et-al-icaps2020}, whose negative result for \bddastar\ the
companion paper answers with the width condition; our contribution here
is orthogonal to both: the operational cost that neither captures.
SetA\ensuremath{^*} and \ghseta\ \citep{jensen-et-al-aij2008} and Lazy
\bddastar\ \citep{torralba-phd2015} vary \emph{when} the partition is
computed but expand the same buckets, so Proposition~\ref{prop-calls}
applies to them as well; the operator-potential planner of
\citet{fiser-et-al-aaai2022,fiser-et-al-aij2024} instantiates \ghseta\
with LP potentials and is the many-valued extreme of our dichotomy.
Closest to pruning-only search are the symbolic perimeter abstractions
of \citet{torralba-et-al-ijcai2016a,torralba-et-al-aij2018}, which also
delay heuristic effort (until blind search stalls; we delay until an
incumbent exists) but use the abstraction to order the frontier;
our variant never partitions and carries the factor-$Wn$ bound.
Upper-bound pruning itself is classical in explicit branch-and-bound;
its symbolic form here is distinguished by the exact slice
characterization (Lemma~\ref{lem-slice}) that makes the guarantee
possible.

\section{Conclusions}

Node counts were the wrong ledger. Counting image operations explains
what a decade of experiments with heuristic symbolic search could not:
$f$-ordering multiplies image calls by the number of heuristic values,
so it pays exactly for few-valued heuristics---merge-and-shrink gains
coverage, potentials lose it---and pruning-only search, which keeps
blind search's image-call profile and a factor-$W$ tighter effort
bound, is confirmed free but gains only what its discards are worth.
Our bidirectional study sharpens the frontier: across four designs,
heuristic pruning never reliably recouped the cost of computing the
heuristic against blind bidirectional search, while seeding an incumbent
from a satisficing planner improved even the blind configuration---the
cheapest useful guidance requires no heuristic at all. The design rules
for future heuristic symbolic search follow directly: order only with
few values, prune with the rest, acquire incumbents early, and treat
heuristic construction cost as a first-class quantity. Turning these
rules into algorithms---coarse structural bucketing with a chosen image
budget, synthesis that minimizes the value count rather than magnitudes,
and guidance from information the search computes anyway---is where we
are headed next.

% ------------------------------------------------------------------------
% Bibliography from the shared aibasel/bib repo (make update-bib).
\bibliographystyle{plainnat}
\bibliography{bib/abbrv,bib/literatur,bib/crossref}
\end{document}